\chapter{Vaqueros: el mismo procedimiento en el papel}

\label{cap:vaqueros}
\section{Trece días de diferencia}

Vaqueros aprobó su Plan de Desarrollo Urbano Ambiental por la \textbf{Ordenanza 474/15}, sancionada en octubre de 2015 y promulgada el \textbf{12 de noviembre} por el Decreto 119/15, trece días antes de que La Caldera aprobara el suyo por la Ordenanza 643. El plan de Vaqueros se elaboró en el marco del Programa de Mejora de la Gestión Municipal BID 1855-OC-AR de la Subsecretaría de Financiamiento provincial; el de La Caldera, con fondos del Consejo Federal de Inversiones (capítulo \ref{cap:pdua}). El 26 de diciembre de 2016, por la \textbf{Ordenanza 507/16} ---promulgada el mismo día por el Decreto 165/16, con firma del intendente Daniel Moreno---, aprobó el «Ordenamiento Normativo Urbanístico» de la ciudad: el Código de Planeamiento Urbano Ambiental (Anexo I), un Código de Edificación que reemplazó al de la Ordenanza 316/2008 (Anexo II) y unos Lineamientos del Área Turismo (Anexo III), vigentes a los diez días corridos de la promulgación. La misma ordenanza subordinó al Código el Reglamento de Fraccionamiento de Tierras de la Ordenanza 437/2014 y condicionó el visado municipal previo que exige el Decreto provincial 1.410/1973 a la verificación de la zonificación. El Código se reformó en 2023 (Ordenanza 607/23) y en enero de 2026 (Ordenanza 648/26, del 26 de enero). El texto que este capítulo analiza es el que publica el Colegio de Arquitectos de Salta con el rótulo «2026», es decir, la versión posterior a esa última reforma. La Caldera, por su parte, reglamentó su plan con un Código de 152 artículos ese mismo año de 2016: los dos municipios se dieron plan en noviembre de 2015 y código en 2016. \textit{Fuentes: el índice oficial de ordenanzas de la Municipalidad de Vaqueros y los textos de las Ordenanzas 469/15, 474/15, 507/16, 535/18, 567/20, 607/23 y 648/26, con sus decretos de promulgación. La 607/23 sustituye sólo los artículos 17 a 19 del Código y su plano de zonificación ---crea las zonas Mixto 1, IC y el polígono RENABAP--- y la 648/26 agrega los distritos M2 Especial, M2 Costanera y PU Ribera; ninguna toca el artículo 14. El texto consolidado que acompaña a la 607/23, de noviembre de 2023, ya lo trae con la misma redacción que el de 2026, palabra por palabra; su origen en la versión de 2016 es, entonces, lo más probable, aunque no está confirmado.} \pendiente{texto original del Anexo I de la Ordenanza 507/16, que la copia obtenida trae sólo como carátula, para confirmar la redacción de 2016 del artículo 14}

Esa simultaneidad es el dato estructural del capítulo. Los dos municipios del departamento se dieron su plan el mismo mes, y Vaqueros --- el municipio más poblado del departamento, segundo componente del aglomerado Gran Salta, donde el mercado de suelo opera desde la resolución de la sucesión Serrey\index{Serrey, Carlos} --- lo tenía aprobado antes del fideicomiso Lares de la Inmaculada, antes de la audiencia pública en el colegio de la congregación, antes del rechazo del amparo vecinal y antes de las manifestaciones por agua potable. En septiembre de 2018, siete meses antes de aquella audiencia, dictó además la Ordenanza 535/18 de «nuevas urbanizaciones», firmada por el presidente del Concejo, Tane Da Souza Correa: un manual de procedimiento en seis pasos ---certificado de uso conforme, prefactibilidad con aprobación del Concejo, certificado de aptitud técnica de Inmuebles, permiso de obra y licencia de comercialización previa garantía, final de obras y matrícula---, con certificado de no inundabilidad de Recursos Hídricos, cesión del diez por ciento para espacio verde, el cinco para el fondo municipal de tierras y el diez para equipamiento, y la exigencia de que la audiencia pública de la Ley 7070 se haga «en instalaciones públicas del Municipio». Rige para los trámites en curso sin proyecto definitivo aprobado y modifica parcialmente la Ordenanza 437. Su visto registra además que el proyecto había sido presentado por el Ejecutivo en 2016 y no aprobado entonces por el Concejo, «trabajada en conjunto entre todos los Municipios del Área Metropolitana», y que los criterios se acordaron en un Comité Metropolitano de Ordenamiento Territorial. Si el trámite de Lares estaba alcanzado por esa ordenanza, y por qué su audiencia se hizo en el colegio de la congregación, es algo que este libro no establece. Lo que distingue a los dos municipios, entonces, no es la fecha sino lo que cada uno escribió, y eso es lo que sigue.

\begin{ficha}{Código de Planeamiento Urbano Ambiental del Municipio de Vaqueros, texto de 2026}
\begin{itemize}[leftmargin=*,nosep]
\item \textbf{Norma que aplica:} PIDUA-VAQ, estructurado en tres escalas: metropolitana, municipal y local.
\item \textbf{Ámbito:} ejido municipal de \textbf{Vaqueros y Lesser}. \textit{La Ley 4365 de 1970 ubicaba las fincas Lesser y Abra de Lesser «del municipio de La Caldera», al oeste de Vaqueros, y la Ley 4987 de 1974 no se refiere a ellas (capítulo \ref{cap:cot}): La Ordenanza 469/15, promulgada el 30 de octubre de 2015 por el Decreto 112/15, fija por coordenadas el ejido urbano de la localidad de Vaqueros y, en su artículo 2º, el «ejido urbano de la localidad de Lesser, jurisdicción del municipio de Vaqueros». Se funda en el artículo 98 de la Ley 1349 y en un motivo fiscal ---zonas con servicios municipales que no tributaban---, y da a Lesser por «comprendido dentro de este Municipio» sin citar ley de límites alguna. Qué norma provincial sostiene esa jurisdicción queda por verificar.}
\item \textbf{Estructura:} 6 capítulos, 19 artículos, 3 anexos de planos.
\item \textbf{Supremacía (art. 3):} el Código de Edificación, el Reglamento de Fraccionamiento de Tierras y toda otra normativa urbanística municipal quedan subordinados.
\item \textbf{Zonas (art. 17):} R, R1, R2, C, UP, UF, \textbf{APHid} (protección hídrica), Rural, Mixto 1, IC, \textbf{Polígono RENABAP}, M2 Especial, M2 Costanera y \textbf{PU Ribera}.
\item \textbf{Altura máxima (art. 18):} 7 metros, excepto 12 en M1 e IC.
\item \textbf{Superficie mínima:} R 450 m²; R1 650 m²; R2 450 m²; M1 800 m²; M2 Especial 10.000 m² con frente mínimo superior a 50 m; IC 1.500 m².
\end{itemize}
\end{ficha}

\section{El artículo 14}

El artículo 14 se titula ``Procesos de Gentrificación'' y establece que la condición periférica metropolitana de Vaqueros no implica su consideración como mera ciudad dormitorio de Salta capital \textbf{ni el desarrollo en su territorio de barrios cerrados no integrados a la vida del municipio}. Dispone además que los usos del suelo, la localización de residencia, la urbanización de parcelas rurales y el accionar de los agentes vinculados al desarrollo inmobiliario no podrán derivar en procesos en los que la población económicamente desfavorecida o vulnerable deba migrar al resultar desplazada por población de mayor nivel adquisitivo, mientras esos sectores se renuevan arquitectónicamente y se valorizan por efecto de la inversión privada.

El artículo 15, inciso d, complementa: promover zonificaciones de densidades medias con lotes pequeños \textbf{con el fin de reducir los procesos de gentrificación} y generar oferta residencial para los pobladores locales.

Un código municipal argentino que nombra la gentrificación, la define técnicamente, la prohíbe y le opone un instrumento de zonificación es una rareza que este libro registra sin ironía. El municipio identificó con precisión el mecanismo que la literatura académica había documentado y que el capítulo \ref{cap:tierra} sostiene como eje: el mismo proceso que trae vecinos nuevos expulsa vecinos viejos.

Y sin embargo el artículo 14 carece de sanción, de autoridad de aplicación designada, de procedimiento de verificación y de indicador que permita constatar cuándo se ha violado. Enuncia una prohibición sin construir el aparato que la vuelve operable.

Una versión anterior de este capítulo leía la comparación con La Caldera como de signo inverso: La Caldera, con el procedimiento completo en su Código y sin aplicarlo; Vaqueros, con el diagnóstico correcto y sin procedimiento. \textbf{Leídas sus ordenanzas, la comparación es de signo igual.} Vaqueros no tiene el procedimiento en su Código de Planeamiento, pero lo tiene en otras dos normas: el Reglamento de Fraccionamiento de Tierras de la Ordenanza 437/14 y el manual de nuevas urbanizaciones de la Ordenanza 535/18. Lo que al artículo 14 le falta ---autoridad, trámite, garantía, sanción--- está allí: certificado de uso conforme, prefactibilidad aprobada por el Concejo, línea de ribera y certificado de no inundabilidad, estudio de impacto con audiencia, garantía de obra, licencia de comercialización y multas por vender sin aprobación. \textbf{Los dos municipios del departamento tienen en el papel el mismo aparato}, casi paso por paso, y la pregunta que este libro hace para uno vale para el otro: si alguien lo aplicó.

\section{La Ordenanza 437 y el anexo que el Boletín no publicó}

El procedimiento de fraccionamiento de Vaqueros existe desde diciembre de 2014. Lo que durante casi una década no existió fue su publicación: el Boletín nunca publicó el reglamento, y el digesto municipal lo ofrece al menos desde febrero de 2024. Este relevamiento lo obtuvo recién al final, por una vía distinta del Boletín.

\begin{cautela}
Lo que el Boletín Oficial publicó de la Ordenanza 437 ---el 15 de enero de 2015--- son los cuatro artículos que aprueban un reglamento; el reglamento mismo no. Durante años, la norma que regula la conversión de suelo rural en urbano en el municipio más presionado del departamento fue de acceso presencial. El texto completo está en el \textbf{Digesto Jurídico Virtual} de la Municipalidad de Vaqueros, una carpeta pública de Google Drive (\texttt{drive.google.com/drive/folders/16YHEfOCD\_j5etW96ChGgO4XtgiRR9Kol}) con más de doscientas treinta ordenanzas y decretos, que el municipio \textbf{ya ofrecía el 2 de febrero de 2024}: ese día su página oficial «recordaba» a vecinos y comerciantes que podían consultarla. De allí lo tomó este libro en septiembre de 2026, y el recordatorio sugiere que el digesto es anterior a esa fecha. Desde 2018, además, la Ley 8126 obliga a los municipios a publicar sus ordenanzas en el Boletín Oficial y en internet (art.\ 79). Es el primero de los cinco casos que el capítulo \ref{cap:metodo} describió, y el único que terminó por resolverse del lado del municipio.

Un dato adicional a verificar: la vicepresidenta del Concejo que sancionó esta ordenanza, publicada en enero de 2015, comparte nombre con la candidata a senadora suplente por el departamento que integró en 2025 la fórmula encabezada por Miguel Calabró. Si se confirma con documento, se trata de una trayectoria de más de una década en los órganos que regulan el suelo del departamento.
\end{cautela}

\subsection*{El texto de la ordenanza, y dónde quedó su anexo}

Se obtuvo el texto de la Ordenanza 437 tal como se publicó en el Boletín Oficial Nº 19.462 del 15 de enero de 2015.

\begin{ficha}{Ordenanza Nº 437 --- Reglamento de Fraccionamiento de Tierras}
Publicada en el Boletín Oficial Nº 19.462 del 15 de enero de 2015 ---orden de publicación 100045820, factura Nº 95, importe \$95---. \textbf{El cuerpo publicado no consigna el número de la ordenanza, la fecha de sanción ni la promulgación}: el número sólo aparece en la leyenda que el Boletín agrega al pie, y la versión «copia fiel» del aviso es idéntica. Como el artículo 2º hace regir la norma desde el día siguiente a su promulgación, \textbf{el texto publicado no permite saber desde cuándo está vigente}. Este libro había registrado la sanción como del 23 de diciembre de 2014; el índice oficial de ordenanzas de la Municipalidad de Vaqueros la consigna como Ordenanza 437/14, del \textbf{29 de diciembre de 2014}, que es la fecha del Decreto 78/14 que la promulga; según su artículo 2º, entonces, rige desde el 30 de diciembre de 2014. La copia oficial del digesto municipal no trae la fecha de sanción. Fundamento: art.\ 21 inc.\ 38 de la Ley 1349. Firman Eloy Sergio García, presidente; María Fernanda Dip Torres, vicepresidenta; y Nelson Damián Gutiérrez, concejal.

\textbf{Objetivos declarados en los considerandos:} crear condiciones normativas para que todo fraccionamiento se realice ``de acuerdo a las mejores formas de utilización y mejoramiento del medio ambiente y desarrollo sustentable''; garantizar ``la compatibilidad ambiental y funcional entre las áreas urbanizadas y a urbanizar''; \textbf{``lograr el máximo aprovechamiento de la infraestructura existente, evitando la habilitación de nuevas áreas sin disponibilidades de extensión de las mismas''}; preservar las áreas de interés natural o paisajístico; y posibilitar la conexión con las áreas existentes mediante un trazado racional de la red vial.

\textbf{Artículo 1º:} ``Apruébase el Reglamento de Fraccionamiento de Tierras en jurisdicción municipal de Vaqueros, \textbf{el que obra en anexo por separado}, que se adjunta y forma parte integrante de la presente''. \textbf{Artículo 2º:} las disposiciones de «este código» rigen desde el día siguiente al de su promulgación. \textbf{Artículo 3º:} deroga toda disposición que se le oponga. \textbf{Artículo 4º:} de forma.

Al pie del aviso, el Boletín Oficial imprime: ``\textbf{El Anexo de la Ordenanza Nº 000437 de la Municipalidad de Vaqueros se encuentra para su consulta en oficinas de esta reparación}'' ---así, «reparación», en el original---. \textbf{La ficha del instrumento en la edición web del Boletín no ofrece enlace «Ver anexo»}, a diferencia de las resoluciones del régimen de tierra fiscal (capítulo \ref{cap:opacidad}).
\end{ficha}

Tres cosas.

\textbf{Primera: los objetivos contradicen lo que ocurrió.} El inciso c se propone ``lograr el máximo aprovechamiento de la infraestructura existente, evitando la habilitación de nuevas áreas sin disponibilidades de extensión de las mismas''. Es decir: no habilitar suelo donde no se pueda llevar servicios. Ese principio, publicado en enero de 2015, convive con un municipio cuya zona alta llevaba entonces al menos ocho años reclamando agua y con clubes de campo a los que el Estado exigía resolver el abastecimiento en forma privada. La ordenanza enuncia exactamente el criterio que el corredor no siguió.

\textbf{Segunda: el anexo es la ordenanza.} El artículo 1º dice que el reglamento ``obra en anexo por separado'' y que ese anexo ``forma parte integrante'' de la norma. Lo publicado en el Boletín, entonces, no es el Reglamento de Fraccionamiento de Tierras: es la ordenanza que lo aprueba. El reglamento mismo --- superficies mínimas, cesiones, obras exigibles, procedimiento --- nunca se publicó allí.

\textbf{Tercera:} según la leyenda, el anexo no quedó en el municipio sino ``en oficinas de esta reparación'' ---así, por «repartición»---, y la repartición que lo dice es \textbf{el Boletín Oficial de la Provincia}, en Avenida Belgrano 1349 de la ciudad de Salta, a una docena de kilómetros del municipio cuya reglamentación contiene. Este libro lo buscó primero en Vaqueros, después allí, y lo encontró finalmente en el digesto que el propio municipio publica en línea.

\subsection*{Qué facultad ejerce esa ordenanza}

La Ordenanza 437 se funda en el artículo 21 inciso 38 de la Ley 1349. Obtenido el texto de esa ley, el inciso dice que es atribución y deber de los Concejos Deliberantes ``dictar planes reguladores para el futuro desarrollo de las ciudades y pueblos y disponer la extensión del ejido urbano, todo de conformidad a lo prescripto por el Capítulo XIII de la presente Ley''.

El reenvío es lo importante, porque ese Capítulo XIII era el piso que la ley orgánica fijaba a los municipios en materia de loteo ---el régimen provincial de aprobación de planos fue, entre 1973 y 2019, el Decreto 1.410/73 (capítulo \ref{cap:loteo})---, y era notablemente delgado.

\begin{ficha}{Capítulo XIII, Ley 1349 (derogada en 2018) --- de las ciudades y pueblos de los municipios}
\begin{itemize}[leftmargin=*,nosep]
\item \textbf{Art.\ 98.} Las Municipalidades y Comisiones Municipales fijan por ordenanza el ejido urbano, \textbf{con aprobación del Poder Ejecutivo} y previo informe de las reparticiones correspondientes.
\item \textbf{Arts.\ 100 y 101.} Fijado el ejido y levantado el plano, se remite copia a Rentas, al Registro de la Propiedad y a Obras Públicas; Rentas clasifica como urbanos los inmuebles comprendidos.
\item \textbf{Art.\ 102.} En casos de fraccionamiento o loteo para la venta dentro del ejido, el propietario \textbf{debe entregar a la Municipalidad un plano de la división practicada} y copia a Rentas. ``En caso de que al hacerse la división se proyecten calles nuevas o pasajes, deberá solicitarse previamente autorización de la autoridad municipal correspondiente.''
\item \textbf{Art.\ 103.} Para la formación de pueblos nuevos, y mientras no exista ley especial sobre la materia, ningún propietario puede subdividir sin previa aprobación del plano y fijación del ejido por la Municipalidad.
\end{itemize}
\end{ficha}

Léase el artículo 102 despacio, porque es el piso normativo municipal sobre el que se construyó buena parte del corredor.

La obligación central del loteador, en la ley orgánica provincial ---vigente hasta 2018---, era \textbf{entregar un plano}. No hay garantía de obra. No hay infraestructura exigible. No hay cesión obligatoria de suelo para equipamiento o espacio verde. No hay estudio de riesgo hídrico. La única autorización previa que la ley pide se dispara por un hecho puramente geométrico: que el fraccionamiento proyecte calles nuevas o pasajes.

De donde se sigue una consecuencia que reordena este capítulo. El Reglamento de Fraccionamiento de Tierras de Vaqueros no reglamenta un régimen exigente: \textbf{lo construye}. Todo lo que exige ---superficies mínimas, cesiones, obras a cargo del desarrollador--- es exigencia que el municipio agrega donde la ley orgánica no la tenía. El Decreto 1.410/73 sí exigía obras antes de la venta, pero ponía su verificación en manos de la propia Municipalidad.

Obtenido el anexo, el régimen construido resulta completo.

\begin{ficha}{Reglamento de Fraccionamiento de Tierras --- anexo de la Ordenanza 437/14 (92 artículos)}
\begin{itemize}[leftmargin=*,nosep]
\item \textbf{Cursos de agua (arts.\ 14, 49 y 62):} en inmuebles afectados por arroyos no se inicia ningún trámite sin la línea de ribera certificada por Recursos Hídricos; no se admiten proyectos que alteren el escurrimiento ni parcelas en canales naturales, cañadas o arroyos.
\item \textbf{Pendiente (art.\ 12):} máxima longitudinal admisible del 12 por ciento.
\item \textbf{Cesiones (art.\ 26):} no menos del 10 por ciento de la superficie de las parcelas para espacios verdes, más un 5 por ciento para uso comunitario.
\item \textbf{Infraestructura (arts.\ 30 a 45):} agua corriente, energía, alumbrado, arbolado, parquización, cordón cuneta y enripiado, a cargo exclusivo del loteador; el agua, con certificado de factibilidad del prestador y plano aprobado por él.
\item \textbf{Ambiente y plazos (art.\ 61):} estudio de impacto ambiental y social conforme la Ley 7070; plan de obras con plazo máximo de cinco años.
\item \textbf{Garantía (arts.\ 68 y 69):} depósito, hipoteca o títulos por el presupuesto de cada obra más un 20 por ciento, reajustable, que asegura la ejecución en los plazos del plan de obras (art.\ 61 inc.\ h).
\item \textbf{Caducidad (art.\ 72):} si las obras no empiezan en un año o no terminan en plazo; si se vendió antes de la aprobación, el municipio ejecuta las obras con la garantía.
\item \textbf{Venta (arts.\ 73, 74 y 84):} prohibido comprometer en venta antes de la aprobación del loteo, cartel obligatorio en el predio y multa por cada parcela vendida sin aprobación.
\end{itemize}
La Ordenanza 535/18 modificó parcialmente sus artículos sobre urbanizaciones, sin decir cuáles, y agregó el manual de seis pasos que se describe al comienzo del capítulo.

\textit{Erratas del original, que se transcriben como están y se aclaran en el apéndice \ref{cap:normativa}: el artículo 26 remite al «artículo 4° Inc.\ a)» para definir la urbanización, que está en el artículo 2º inciso a; y el artículo 68 remite, para los plazos que la garantía debe cubrir, al «art.\ 71° inc.\ f» en la copia firmada ---«61° inc.\ f» en la versión transcripta del mismo digesto---, cuando los plazos están en el artículo 61 inciso h.}
\end{ficha}

Lo que faltaba, entonces, no era la regla: era su texto. Y con el texto a la vista, lo que el apartado siguiente describe deja de poder atribuirse a la ausencia de norma.


\section{Y lo que ese reglamento debía impedir, ocurrió}

El principio de la Ordenanza 437 ---no habilitar suelo donde no se pueda llevar servicios--- está publicado en
enero de 2015.

Tres meses antes, en octubre de 2014, se había celebrado en el Centro Deportivo Municipal de
Vaqueros la audiencia pública ambiental por la Urbanización Cerro de Buena Vista. Un año
antes, en enero de 2014, el Ministerio provincial había aprobado el plano del club de campo
Las Vertientes exigiendo un reglamento de condominio que contemplara el \textbf{servicio de
agua en forma privada}. Y siete años después, en julio de 2021, los vecinos de la zona alta
del mismo municipio se manifestaban tras más de quince años de reclamos por agua potable.

\textbf{La ordenanza enuncia exactamente el criterio que el municipio no aplicó}, y lo enuncia
un año después de que el Estado provincial homologara, sobre ese territorio, un
abastecimiento privado.

No se sigue de ahí ninguna imputación: el reglamento que fija cómo se aplica ese criterio es
el anexo que el Boletín no publicó, y que este libro obtuvo al final del digesto municipal:
exige agua corriente, línea de ribera, cesiones y garantía antes de vender. Se sigue una
constatación de forma, que es la del capítulo entero: \textbf{en Vaqueros el diagnóstico y el
procedimiento están escritos; lo que no consta es que alguien los haya aplicado.}

\section{Las zonas que registran el conflicto}

Tres zonas del artículo 17 son el reconocimiento normativo de conflictos documentados. \textbf{APHid}, área de protección hídrica, cuya definición completa es de tres líneas: debe preverse futura infraestructura de defensas; a futuro será UP. \textbf{PU Ribera}, de utilización prioritaria como espacio abierto, concordante con el artículo 15 inciso b, que ordena recuperar los cauces y márgenes de los arroyos, evitar construcciones sobre líneas de ribera y jerarquizar los caminos de sirga como espacios verdes públicos. Y \textbf{Polígono RENABAP}, uso residencial de densidad alta: el reconocimiento formal de que existe en Vaqueros al menos un barrio popular registrado.

\subsection*{De qué tamaño es el conflicto que esas tres líneas reconocen}

La fuente quedó identificada: es la presentación \textit{Línea de Ribera y Ordenamiento Territorial}, del geólogo Jorge G.\ Torres\index{Torres, Jorge G.}, de la Secretaría de Recursos Hídricos de Salta, publicada por el Consejo Hídrico Federal. Es documento oficial y su autor es hoy Director General de Cuencas y Política Hídrica de la provincia.

Con una salvedad que corresponde hacer. La presentación incluye una comparación fotográfica del cauce del río Vaqueros rotulada ``\textbf{Año 2019}'' y ``\textbf{Año 2021}'', pero \textbf{las cifras de 229 y 89 metros no figuran en su texto}: provienen de un resumen intermedio de ese documento. Hasta poder leerlas en el original --- o medirlas sobre las mismas imágenes ---, este libro \textbf{no publica esos números}. Lo que en cambio está en el documento, y en sus conclusiones, es la afirmación cualitativa: los emprendedores inmobiliarios ``solicitan la modificación de sectores para ganar espacios y \textbf{avanzar hacia el dominio fluvial}''.

La presentación describe dos mecanismos que operan sobre las riberas del Valle de Lerma:

\begin{itemize}[leftmargin=*]
\item \textbf{Corrimiento de la línea de ribera a pedido de parte.} ``En este punto también es en algunos casos conflictivo, ya que algunos emprendedores quieren la modificación de la Línea de Ribera, para poder cumplir con el 5\% requerido por el municipio.'' Y se menciona ``un caso reciente, y aún no resuelto'': una urbanización privada ``donde originalmente el espacio verde fue ocupado por la construcción de dúplex''. El documento no identifica esa urbanización ni su municipio.
\item \textbf{Extracción de áridos sin control.} La extracción por debajo de un metro y medio de profundidad altera los canales subálveos, y la acumulación de descartes dentro del cauce reduce su ancho útil y genera anegamientos.
\end{itemize}

\begin{cautela}
El primer mecanismo merece detenerse, porque invierte el sentido de un instrumento de protección.

La cesión de espacio verde es una exigencia municipal destinada a que el loteo entregue superficie a la comunidad. La línea de ribera es una determinación técnica destinada a fijar hasta dónde llega el agua. Si un desarrollador puede pedir que la segunda se corra para cumplir la primera, entonces \textbf{el porcentaje de espacio verde deja de ser una cesión y pasa a ser un cálculo}: no se entrega suelo aprovechable, se reclasifica suelo ribereño hasta que la cuenta cierre. El espacio verde exigido termina siendo, exactamente, la franja que el río reclamará.

La puerta existe y ahora se conoce su marco completo. El artículo 126 del Código de Aguas faculta a la Autoridad a rectificar ``cuando el cambio de las circunstancias lo haga necesario'', y el artículo 1º del Anexo I del Decreto 1989/02 admite que la determinación se realice \textbf{a solicitud de partes interesadas}. El pedido de parte es, entonces, un cauce previsto por la norma, no una anomalía.

Que la puerta exista sigue sin decir quién la usa ni para qué, y este libro no lo afirma respecto de Vaqueros ni de ningún loteo. La diferencia con lo que este capítulo podía decir antes es que ya no se trata de un mecanismo descripto por una fuente incierta: lo describe un funcionario provincial del área, en un documento oficial, y el capítulo \ref{cap:ribera} registra una rectificación concreta y fechada en el departamento --- la Resolución 108/2025 sobre el río La Caldera.
\end{cautela}

\subsection*{Seis años de crecidas sobre el mismo río}

La reducción del cauce que un resumen intermedio atribuye al río Vaqueros ---con cifras que este libro no da por verificadas--- tiene, se confirme o no, una serie de correlatos públicos que sí pueden fecharse.

\begin{ficha}{Serie de prensa sobre el río Vaqueros, 2020--2026}
\begin{itemize}[leftmargin=*,nosep]
\item \textbf{Abril de 2020.} Las lluvias de principios de mes ``voltearon las defensas a la altura del \textbf{barrio 15 de Septiembre}''. Varias casas en riesgo de desmoronarse; sus habitantes ``no tienen dónde ir'' y piden asistencia estatal. La nota se clasifica bajo desigualdad y pobreza y se etiqueta a la Municipalidad de Salta.
\item \textbf{Obras de encauzamiento} ejecutadas por la Municipalidad de Salta, presentadas como beneficio para \textbf{1.700 familias}.
\item Vecinos de las márgenes \textbf{en alerta por el avance del cauce}; una crecida \textbf{destruye un muro de contención}.
\item \textbf{Enero de 2026.} Rescate de \textbf{más de treinta personas} por crecidas de ríos en La Caldera, Vaqueros y Campo Quijano. La crecida se lleva el terraplén y una casilla con herramientas de obra pública.
\item \textbf{19 de enero de 2026.} La crecida ``arrasó una parte clave de la obra del \textbf{nuevo puente} que unirá Vaqueros con Salta''.
\item El gobernador supervisa las obras de \textbf{puentes sobre el río Vaqueros y Circunvalación}, con asistencia financiera provincial.
\end{itemize}
\textit{Advertencia: de esta serie se leyó completo únicamente el primer ítem, y de él sólo el copete. Los restantes se consignan por sus titulares. Antes de citarlos hay que leer cada nota y verificar fecha y contenido.}
\end{ficha}

Conviene separar tres asuntos que la serie permite decir con la prudencia del caso.

\textbf{La primera es que el río Vaqueros es un límite, no un accidente interno.} Separa el municipio de Vaqueros --- departamento La Caldera --- de la ciudad de Salta. Por eso el barrio afectado en 2020 aparece bajo jurisdicción capitalina y las obras de encauzamiento las ejecuta la Municipalidad de Salta. \textbf{El cauce cuyo estrechamiento importa a este capítulo tiene dos municipios en sus dos orillas y ninguno con competencia sobre el conjunto.} El artículo 15 inciso b del Código de Vaqueros ordena recuperar cauces y márgenes; puede hacerlo sobre una sola margen.

\textbf{La segunda es que la presión no viene sólo del loteo.} Un puente nuevo entre Vaqueros y Salta, obras de circunvalación y encauzamientos son infraestructura que reduce la fricción entre el corredor y la capital, y por lo tanto valoriza el suelo. Este libro documentó que la conversión precede a los servicios; la serie de prensa anterior muestra el otro extremo del ciclo, donde el servicio llega y consolida lo convertido.

\textbf{Y la tercera, que es la más incómoda:} como muestra esa misma serie, la crecida de enero de 2026 dañó la obra construida para ordenar el río. Una infraestructura pensada para contener el cauce fue alcanzada por el cauce antes de terminarse. Es el mismo patrón que el capítulo \ref{cap:ribera} documenta aguas arriba, y merece la misma pregunta: con qué línea de ribera se proyectó.


\subsection*{Y hay una obra mayor, con audiencia pública convocada}

En noviembre de 2023 la Dirección de Vialidad de Salta convocó a audiencia pública sobre el
Estudio de Impacto Ambiental y Social de una obra que este libro no había registrado.

\begin{ficha}{Boletín Oficial, 14, 15 y 16 de noviembre de 2023}
«Estudio de Impacto Ambiental y Social --- Obra: \textbf{Pavimentación Ruta Nacional Nº 9,
tramo Salta -- La Caldera}. \textbf{Sección I: Rotonda Avda.\ Bolivia -- Puente Río Caldera.
Sección II: Puente Río Vaqueros -- Puente Río Caldera}. Expte.\ Nº
\textbf{0110033-243466/2021-181}. Dirección de Vialidad de Salta.»

Marco invocado: \textbf{artículo 49 de la Ley 7070} y \textbf{Resolución DVS 457/2007}, sobre
el Reglamento de Audiencia Pública de Proyectos Viales.

\textbf{Audiencia: 27 de noviembre de 2023, a las 9, en el Centro Integrador Comunitario de la
Municipalidad de Vaqueros.} Plazo para presentarse como parte: hasta las 13 del 22 de noviembre.
El expediente podía consultarse \textbf{en el Consejo Técnico de Vialidad, España 721, ciudad
de Salta, de 9 a 13}.
\end{ficha}

\textbf{Es la obra de mayor alcance sobre el corredor que este libro registra.} Va de la rotonda
de avenida Bolivia ---el borde de la capital--- hasta el puente del río La Caldera, y tiene una
segunda sección entre los dos puentes. Es, exactamente, la reducción de fricción que el párrafo
anterior describe como mecanismo de valorización, ejecutada sobre la traza completa.

\begin{cautela}
\textbf{Y su convocatoria repite, punto por punto, la forma que este libro viene describiendo.}

\textbf{La audiencia sobre una obra que termina en La Caldera se convocó en Vaqueros.} Quien vive
en el pueblo tenía que viajar al otro municipio, un lunes a las nueve de la mañana.

\textbf{El plazo para constituirse en parte fue de ocho días}, contados desde la primera
publicación. Para una obra vial de dos secciones sobre la ruta nacional que atraviesa el
departamento.

\textbf{Y el expediente sólo podía verse en la capital}, en España 721, de nueve a una. Es la
misma restricción que el capítulo \ref{cap:ribera} documenta para las coordenadas de línea de
ribera, sobre otro organismo y nueve años antes: \textbf{el acto se publica, el plazo corre, y
el documento que permitiría opinar está a veinte kilómetros y en horario de trabajo}.

\textbf{Este libro no sabe qué pasó en esa audiencia}, ni quiénes se presentaron, ni si alguien
de La Caldera lo hizo, ni en qué estado está la obra hoy. El apéndice \ref{cap:pedidos} lo pide.
\end{cautela}

\subsection*{Y en 2026 la cobertura de la obra vial dice en voz alta lo que el libro venía documentando}

Casi tres años después de aquella audiencia, el director de Vialidad Nacional en la provincia
describe públicamente el estado del corredor y su límite.

\begin{ficha}{Declaraciones publicadas el 21 de agosto de 2026 en \textit{Construar}}
Los trabajos en curso ---despeje de banquinas, refuerzo de sectores críticos y limpieza de obras
de drenaje--- avanzan \textbf{desde el límite provincial hasta el puente Wierna} y se esperaba
concluirlos en septiembre.

Sobre el estado de la calzada: «\textit{Hoy tenemos una calzada de \textbf{5 metros, 5 metros y
medio}. Cuando se construyó, hace más de 40 o 50 años, seguramente cumplió su cometido. Hoy está
claro que ha quedado en déficit para el volumen de tráfico que allí circula}».

Sobre lo que impide ampliarla habla la nota, en su propia voz ---sin comillas ni atribución
expresa al funcionario---: «\textit{El crecimiento poblacional de localidades como
Vaqueros y La Caldera, sumado al \textbf{desarrollo de emprendimientos privados sobre las
laderas ---muchos de ellos sin autorizaciones---, complica cualquier ampliación futura del
corredor}}».

\textit{Fuente:} \textit{Construar.com.ar}, «Ruta Nacional 9 en Salta: trabajos de mejora y
debate por el angosto ancho de calzada», 21/08/2026,
\url{https://www.construar.com.ar/2026/08/ruta-nacional-9-en-salta-trabajos-de-mejora-y-debate-por-el-angosto-ancho-de-calzada/}
(consultada el 18/09/2026).
\end{ficha}

\begin{cautela}
\textbf{Esa última frase es lo más importante que este capítulo incorpora, y conviene medir
exactamente su alcance.}

\textbf{Qué es.} Una afirmación de un medio especializado, en una nota que recoge las
declaraciones del titular provincial de un organismo nacional. \textbf{La frase no está entre
comillas}: puede ser paráfrasis de lo que el director dijo o redacción del medio, y la nota no
permite saber cuál de las dos. \textbf{No es un expediente, no individualiza ningún
emprendimiento y no constituye una imputación.} Este libro no puede decir cuáles son, ni
cuántos, ni sobre qué inmuebles.

\textbf{Qué agrega.} La cobertura de un \textbf{tercer organismo} ---después del municipal y del
hídrico--- registra que hay desarrollos privados sobre las laderas del corredor y que muchos carecen de
autorización. El capítulo \ref{cap:loteo} documenta esos desarrollos uno por uno con su
expediente; el capítulo \ref{cap:resistencias} recoge que la autoridad minera informó en 2022 la
existencia de loteos dentro de líneas de ribera delimitadas. Si la frase es del director,
Vialidad Nacional lo dice ahora desde otra jurisdicción y por otra razón; si es del medio, lo
dice la prensa especializada que sigue la obra, y el pedido del apéndice \ref{cap:pedidos}
busca el documento que lo respalde.

\textbf{Y agrega una consecuencia que el libro no tenía.} Hasta acá el problema de los
fraccionamientos era de dominio, de agua y de riesgo. Ahora es también \textbf{de
infraestructura pública}: lo construido sin autorización \textbf{limita lo que el Estado puede
hacer después}. Una calzada de cinco metros y medio, proyectada hace medio siglo, no puede
ensancharse porque las laderas por donde debería crecer ya están ocupadas.

\textbf{Puesto al lado del capítulo \ref{cap:cot}, el contraste es el del libro entero.} En 2014
el Estado fijaba para una ruta provincial de este mismo departamento un ancho de zona de camino
de 14,70 metros declarados y sesenta por ley. En 2026, sobre la ruta nacional que lo atraviesa,
la calzada real es de cinco y medio, según el propio organismo, y ampliarla es, según la nota
que lo entrevista, complicado.
\end{cautela}




La zona APHid describe, sin proponérselo, la lógica del departamento: se protege un área hídrica declarando que en el futuro tendrá defensas y que después será urbanizable. La protección se enuncia como etapa transitoria hacia el uso urbano.

En La Caldera la ribera del río es zona no urbanizable de manera permanente. En Vaqueros el área de protección hídrica es, por definición normativa, urbanización parque diferida. Dos municipios del mismo departamento, sobre la misma cuenca, con criterios opuestos sobre el mismo tipo de suelo.

\section{Lares de la Inmaculada}

De los casos relevados, éste articula en un solo expediente la propiedad eclesial de la tierra, la operación inmobiliaria de escala y la membresía cruzada de los cuadros jurídicos.

\begin{ficha}{Fideicomiso Lares de la Inmaculada}
\begin{itemize}[leftmargin=*,nosep]
\item En septiembre de 2017 se constituyó el fideicomiso. La Congregación de los Hijos de la Inmaculada Concepción transfirió veintitrés hectáreas por esa vía, sobre la matrícula 3012 del departamento La Caldera.
\item El rol de fiduciario fue confiado a un arquitecto de apellido Lobato; por el mismo contrato, la desarrolladora MDay quedó al frente de la urbanización.
\item La cobertura periodística atribuye al proyecto veintisiete hectáreas ---cuatro más que las transferidas, diferencia que no se explica--- y doscientas veintisiete parcelas, en el centro de Vaqueros.
\item Finalidad declarada: destinar los recursos a la ampliación del Colegio San Cayetano, que funciona en el mismo predio, y duplicar su matrícula.
\end{itemize}
\end{ficha}

La audiencia pública prevista en el procedimiento de habilitación ---convocada por el Programa de Audiencias Públicas del Ministerio de Producción, Trabajo y Desarrollo Sustentable como «La Inmaculada Urbanización», expediente 0090227-118038/2018-0, con avisos en el Boletín Oficial del 22, 25 y 26 de marzo de 2019--- se realizó en abril de 2019, presidida por una funcionaria de la Secretaría de Ambiente provincial, \textbf{en el salón de usos múltiples del Colegio San Cayetano} --- esto es, en un inmueble de la propia congregación impulsora. Según la cobertura de \textit{Cuarto Poder}, desde la institución se invitó a padres y madres de alumnos a presentarse, pese a que la instancia debía realizarse en lugar neutral y en pie de igualdad. Dos vecinos habían impugnado previamente la audiencia objetando la falta de dictamen previo del municipio; el recurso no prosperó.

Los vecinos plantearon que no había habido un estudio serio de impacto ambiental, que temían contaminación por la cantidad de familias a instalarse en una zona sin cloacas, y que las viviendas más próximas quedarían entre el loteo y un arroyo que desemboca en el río Vaqueros.

La sede de la audiencia no es un detalle de logística. El procedimiento existe para producir deliberación en condiciones de simetría entre proponente y afectados. Celebrarla en el establecimiento del proponente, con convocatoria dirigida a su comunidad educativa, no viola necesariamente ninguna norma expresa --- y por eso mismo es el ejemplo más limpio de lo que el libro anterior llamó tecnología de coacción de primer grado: la que opera sin infringir nada.

Media docena de vecinos, patrocinados por el abogado Luis Segovia, interpuso amparo contra la Municipalidad de Vaqueros para que el proyecto cumpliera las previsiones de planificación urbana, ordenamiento territorial y ecológico; la Congregación y la desarrolladora fueron admitidas como partes, y se cuestionó la preventa de las doscientas veintisiete parcelas prevista en el contrato. El abogado de la Congregación sostuvo que lo planteado era un interés privado y no público; el de la desarrolladora, que ciento cincuenta familias de la parte baja resultarían beneficiadas, sobre todo en materia de desagües. En noviembre de 2019 la Sala III de la Cámara de Apelaciones en lo Civil y Comercial rechazó el amparo.

Una nota de \textit{Cuarto Poder} de junio de 2026, en cobertura sobre litigios de consumidores contra el grupo desarrollador, consigna que el abogado que lo representó en juicios como el de Lares de la Inmaculada se desempeñaba, en el período en que la intendencia capitalina favorecía otra urbanización del mismo grupo, como abogado de la Municipalidad de Salta.

Si el dato se confirma con documentación --- designación en boletín oficial por un lado, presentaciones judiciales por otro --- constituiría la primera verificación documentada de H.23 en este territorio: el mismo cuadro jurídico operando en el vértice público que regula y en el privado que desarrolla, en municipios distintos del mismo aglomerado y en el mismo período. No es ilegal ejercer la profesión privada y un cargo municipal en jurisdicciones diferentes. Es, sin embargo, la configuración que el libro anterior describió como matriz del dispositivo. Hasta esa verificación, es prensa.


